\section{Introduction}
\label{sec:intro}

Three complaints about contemporary language models recur often enough to be
folklore. Speculative decoding \cite{leviathan2023,chen2023} accelerates
inference by having a second, smaller network guess ahead --- why can a model not
guess ahead using its own early layers? Reasoning is elicited by having models
emit long chains of intermediate tokens --- does that not indicate a model that
never learned to reason internally? And decoder-only transformers, whose
key-value cache is carried forward and updated at every step, are recurrent
networks with an unusually large state --- so why not build the recurrent network
directly?

Each complaint is usually argued on its own. We observe that they are the same
question asked along three different axes: \emph{where is a model permitted to
spend serial computation, and what carries state while it does so?}
Speculative decoding is about the \emph{depth} axis, where the residual stream
already contains a usable intermediate answer that is discarded. Chain of
thought is about the \emph{time} axis, where serial computation is routed
through the vocabulary. The recurrent-form question is about the \emph{state}
axis, where an unbounded cache substitutes for a learned compression.

Our contribution is not a new architecture but a precise accounting of what each
complaint does and does not license. We find the three fare very differently
under formalisation, and in two cases the honest conclusion is not the one the
complaint anticipates.

\paragraph{Depth.} Speculative sampling is exact for an \emph{arbitrary} draft
distribution (Theorem~\ref{thm:exact}), so correctness never justified a separate
draft network; the question is entirely one of acceptance against cost. We give a
cost model interpolating between memory- and compute-bound decoding
(Definition~\ref{def:cost}) under which early-exit self-drafting strictly
dominates any external draft of comparable relative cost, in every regime
(Theorem~\ref{thm:dominance}), and lower-bound the acceptance rate by the
residual tail energy (Theorem~\ref{thm:accept}). We also identify a cost term
that the informal argument omits and that materially changes the conclusion for
small models: every draft step must stream the unembedding, which caps the
attainable speedup (Corollary~\ref{cor:ceiling}) unless the draft head is
shortlisted (Corollary~\ref{cor:shortlist}).

\paragraph{Time.} The strong form of the complaint is false: under
$\mathrm{TC}^0 \ne \mathrm{P}$, serial steps buy computational power that fixed
depth cannot recover (Theorem~\ref{thm:tc0}). The comparison between media is a
genuine trade-off, and it runs \emph{opposite} to the usual claim. A
chain-of-thought step writes only $\log_2 V \approx 17$ bits, but those writes
accumulate; a fixed-size latent loop writes far more per step but overwrites. Past
a computable crossover --- about $238$ steps for a $1024$-dimensional model ---
the token stream is the larger memory (Corollary~\ref{cor:crossover}). The
defensible version of the complaint is therefore that latent reasoning should
\emph{append} rather than overwrite.

\paragraph{State.} The recurrent form is real but nearly vacuous as usually
stated; its content is incrementality, not existence (Theorem~\ref{thm:kvform}).
A constrained dense recurrence does coincide exactly with the HMM forward
algorithm (Theorem~\ref{thm:hmm}). Our main negative result is that this is
\emph{not} a statement about gated linear attention: we prove that no injective
linear embedding carries the forward algorithm into a diagonally-gated
linear-attention recurrence (Theorem~\ref{thm:noembed}), by a support-monotonicity
argument and by simultaneous diagonalisability. Since diagonality is precisely
what makes these architectures parallelisable, ``constrain a gated RNN and obtain
HMM-like decoding'' is false as stated.

\paragraph{Falsifiability.} Every claim above is either proved, refuted, or
labelled open. Section~\ref{sec:unified} gives an experimental protocol at
$\le$1B parameters in which each thesis has a pre-registered kill criterion, and
two of the four experiments require no training.
